\section{Recovery through the Hopf fibration}
\label{sec:hopf-transverse-recovery}

Positive lower volume on a closed base can force recovery even when the ambient manifold has unbounded upper mass. We now take
\[
 X=(\C^3\setminus\{0\})/\langle z\mapsto\mathrm e^{-1}z\rangle,
 \qquad B=\mathbb P^2,\qquad \pi\colon X\to B,
\]
and use the geometry of \cref{chap:hopf}. Write $\kappa=\omega_{\FS}$, with $\int_B\kappa^2=1$, and put
\[
 s=\log|z|^2,\qquad
 \alpha=\pi^*\kappa=\ddc s,\qquad
 \beta=\mathrm ds\wedge\dc s,\qquad h=\alpha+\beta.
\]
The derivatives of $s$ descend to $X$. Scalar multiplication gives an action of the elliptic group $E=\C/(\mathbb Z+2\pi\mathrm i\mathbb Z)$, transitive on each fiber of $\pi$ and preserving $h$. By \cref{thm:hopf-unbounded}, $h$ has unbounded upper mass. Invariant potentials nevertheless reduce to potentials on $B$ because the horizontal part of $h$ is the pullback of $\kappa$.

\Needspace{14\baselineskip}
\begin{proposition}\label{prop:hopf-fiber-ideal-envelope}\index{Hopf manifold!fiber ideal envelope}
For every $b>0$ and $v\in\PSH(B,b\kappa)$,
\[
 P_{bh}[\pi^*v]_{\mathcal I}
       =\pi^*P_{b\kappa}[v]_{\mathcal I}.
\]
The envelope on the left uses all ambient competitors.
\end{proposition}
\begin{proof}
First, an $E$-invariant $q\in\PSH(X,bh)$ descends to an upper semicontinuous function $w$ on $B$, since $\pi$ is proper and $q$ is constant on its fibers. In a product chart, Fubini gives $w\in L^1_{\mathrm{loc}}$, and
\[
 bh+\ddc q=\pi^*(b\kappa+\ddc w)+b\beta.
\]
Lift a smooth base vector to a vector annihilated by $\partial s$. Its pairing with $\beta$ is zero. Testing the positive current above against these lifts and a nonnegative smooth fiber density of integral one proves $b\kappa+\ddc w\geq0$ distributionally. A local potential of $b\kappa$ gives its canonical plurisubharmonic representative. Its pullback agrees almost everywhere with $q$, hence everywhere by \cref{lem:qpsh-almost-everywhere-order}. Thus $w\in\PSH(B,b\kappa)$. Conversely, every such $w$ lifts to a $bh$-plurisubharmonic function by the displayed identity.

Now set $q=P_{bh}[\pi^*v]_{\mathcal I}$. Translation by any $g\in E$ preserves its complete competitor family: it preserves the background, the obstacle zero, and the ideals of the invariant input. The envelope is their greatest member, so $q\circ g=q$. Descend $q=\pi^*w$ as above. By \cref{lem:submersion-multiplier-pullback}, its ideal equivalence to $\pi^*v$ implies $w\sim_{\mathcal I}v$: a germ pulled back from the base is weighted-integrable in a product chart if and only if the base germ is, by Fubini. Hence $w$ is a base competitor. Every base competitor also lifts to an ambient competitor by the same lemma and the current identity. Taking the greatest members gives equality.
\end{proof}

For an envelope system $\boldsymbol\Gamma\in\mathscr S_{a,c}(h)$, we call the components $\Gamma^{\epsilon h}$, $\epsilon>0$, its \emph{scalar components}. If they are invariant under $E$, they alone determine an envelope system on the base, including its components over nonclosed perturbations of $\kappa$.

\Needspace{17\baselineskip}
\begin{lemma}\label{lem:hopf-transverse-system}
Let $\boldsymbol\Gamma\in\mathscr S_{a,c}(h)$, and suppose that $\Gamma_\tau^{\epsilon h}$ is $E$-invariant for every $\epsilon>0$ and $\tau<c$. There is a unique envelope system $\boldsymbol\Psi\in\mathscr S_{a,c}(\kappa)$ such that
\[
 \Gamma_\tau^{\epsilon h}=\pi^*\Psi_\tau^{\epsilon\kappa}
                 \quad(\epsilon>0,\ \tau<c).
\]
Writing $\Gamma=\mathsf R\boldsymbol\Gamma$ and $\psi=\mathsf R\boldsymbol\Psi$, we have $\Gamma_\tau=\pi^*\psi_\tau$ for every $\tau<c$.
\end{lemma}
\begin{proof}
The descent in \cref{prop:hopf-fiber-ideal-envelope} gives
$\Gamma_\tau^{\epsilon h}=\pi^*\phi_\tau^\epsilon$ with $\phi_\tau^\epsilon\in\PSH(B,(1+\epsilon)\kappa)$. Normalization, decrease, concavity and constancy for $\tau\leq a$ descend pointwise. The fiber-envelope identity descends both ideal-model fixedness and scalar compatibility:
\[
 \phi_\tau^{\epsilon_2}
       =P_{(1+\epsilon_2)\kappa}
             [\phi_\tau^{\epsilon_1}]_{\mathcal I}
                 \quad(\epsilon_2\geq\epsilon_1>0).
\]
In particular, the slices $\phi_\tau^\epsilon$, $\epsilon>0$, have the same multiplier ideals at every weight.

For a Hermitian form $\xi>0$ on $B$, choose $\delta>0$ with $\delta\kappa\leq\xi$, fixed for the whole parameter interval, and define
\[
 \Psi_\tau^\xi=P_{\kappa+\xi}
                       [\phi_\tau^\delta]_{\mathcal I}.
\]
The definition is independent of $\delta$ because the scalar inputs have the same ideals. It reproduces the scalar components, and \cref{prop:hermitian-i-curve-transition} proves that each component is a curve in $\mathscr T_{a,c}(\kappa+\xi)$. Projection preserves the ideals, so comparing the complete competitor families proves compatibility for $\xi\geq\xi'$. Thus $\boldsymbol\Psi$ is an envelope system over $\kappa$. Every possible extension is forced by this formula, which proves uniqueness.

The retraction theorem applies to both systems. Taking their decreasing limits along $h/j$ and $\kappa/j$ in the scalar identity gives $\Gamma_\tau=\pi^*\psi_\tau$. In particular the limits, including their values at poles, agree. For any other Hermitian form $k$ on $B$, constants $0<l\leq r$ with $l\kappa\leq k\leq r\kappa$ give
\[
 \phi_\tau^{l/j}\leq\Psi_\tau^{k/j}\leq\phi_\tau^{r/j}.
\]
Both outer sequences decrease to $\psi_\tau$. This also identifies the transverse limit directly, without assuming that it keeps the ideals of the slices $\phi_\tau^\epsilon$.
\end{proof}

\Needspace{15\baselineskip}
\begin{theorem}\label{thm:hopf-transverse-recovery}\index{Envelope system!transverse Hopf recovery}
Let $\boldsymbol\Gamma\in\mathscr S_{a,c}(h)$ have $E$-invariant scalar components as in \cref{lem:hopf-transverse-system}. Let $\psi=\mathsf R\boldsymbol\Psi$ be the transverse limiting curve. If
\[
 \underline{\operatorname{vol}}_\kappa(\psi_a)>0,
\]
then for every $\tau<c$ and every Hermitian form $\eta>0$ on $X$,
\[
 \Gamma_\tau^\eta=P_{h+\eta}[\pi^*\psi_\tau]_{\mathcal I}.
\]
Thus $\boldsymbol\Gamma=\mathsf E(\mathsf R\boldsymbol\Gamma)$.
\end{theorem}
\begin{proof}
The closed K\"ahler background $\kappa$ is big and $B$ has the bounded mass property. Apply \cref{thm:positive-volume-system-recovery} to the full transverse system. Its initial-volume hypothesis is exactly the displayed one, and it gives
\[
 \Psi_\tau^{\epsilon\kappa}
       =P_{(1+\epsilon)\kappa}[\psi_\tau]_{\mathcal I}
                    \quad(\epsilon>0,\ \tau<c).
\]
The fiber-envelope identity yields
$\Gamma_\tau^{\epsilon h}=P_{(1+\epsilon)h}[\pi^*\psi_\tau]_{\mathcal I}$.
For a general $\eta>0$, choose $\epsilon h\leq\eta$. Compatibility and the nested-background identity of \cref{prop:hermitian-i-curve-transition} then give
\[
 \Gamma_\tau^\eta
 =P_{h+\eta}[\Gamma_\tau^{\epsilon h}]_{\mathcal I}
 =P_{h+\eta}[\pi^*\psi_\tau]_{\mathcal I}.
\]
\end{proof}

The positivity condition is a lower ordinary-type volume on the closed base. It does not use an ambient upper-volume deficit or identify the ideals of the limit in advance. The conclusion holds at every finite $\tau<c$; it gives no uniform positive-volume bound as $\tau$ approaches $c$.

\subsection{A nonanalytic family}

The countable hyperplane series of \cref{countable-divisorial:profile,countable-divisorial:positive-mass} supplies an explicit family satisfying the transverse condition. Choose infinitely many distinct hyperplanes $H_j\subset B$, positive coefficients $q_j$ with $\sum_jq_j=1$, and their standard logarithms $\ell_j\leq0$. Put
\[
 u_0=\sum_{j\geq1}q_j\ell_j
                   -\sup_B\sum_{j\geq1}q_j\ell_j.
\]
Then $u_0\leq0$ is $\kappa$-plurisubharmonic and $\kappa+\ddc u_0=\sum_jq_j[H_j]$.

\Needspace{18\baselineskip}
\begin{example}\label{ex:hopf-nonanalytic-recovery}\index{Hopf manifold!nonanalytic recoverable system}
Fix $0<\lambda_0<\lambda_*<1$ and $r>0$ with $r(c-a)<\lambda_*-\lambda_0$. For $\tau<c$, set
\[
 \lambda(\tau)=\lambda_0+r\max(\tau-a,0),\qquad
 u_\tau=\lambda(\tau)u_0,
\]
and, for every Hermitian form $\eta>0$ on $X$, define
\[
 \Gamma_\tau^\eta=P_{h+\eta}[\pi^*u_\tau]_{\mathcal I}.
\]
These components form an envelope system with $E$-invariant scalar slices. Its transverse limiting curve is
$\psi_\tau=P_\kappa[u_\tau]_{\mathcal I}$, and
\[
 \underline{\operatorname{vol}}_\kappa(\psi_\tau)
 =\overline{\operatorname{vol}}_\kappa(\psi_\tau)
 =(1-\lambda(\tau))^2>0.
\]
The system is recovered from its limit. Its ideal data cannot be represented by analytic singularities.
\end{example}
\begin{proof}
The function $\lambda$ is positive, increasing and convex, and $u_0\leq0$, so $u_\tau$ decreases, is pointwise concave and is constant for $\tau\leq a$. Moreover,
\[
 \kappa+\ddc u_\tau
       =(1-\lambda(\tau))\kappa
                         +\lambda(\tau)\sum_jq_j[H_j]>0,
 \qquad
 h+\ddc\pi^*u_\tau\geq(1-\lambda_*)h.
\]
Set $q_\tau=P_\kappa[u_\tau]_{\mathcal I}$. The order, scaling and addition inequalities of \cref{prop:i-envelope-operations} show that $q_\tau$ remains decreasing and concave; it is constant on the same left half-line. Thus $q\in\mathscr T_{a,c}(\kappa)$. Apply the same argument upstairs to $Q_\tau=P_h[\pi^*u_\tau]_{\mathcal I}$. The fiber-envelope identity gives $Q_\tau=\pi^*q_\tau$, and the nested-background identity identifies the displayed family with $\mathsf E(Q)$. It is therefore an envelope system, its limit is $Q$ by $\mathsf R\mathsf E=\mathrm{id}$, and its scalar slices are invariant.

Apply \cref{countable-divisorial:positive-mass} with background $\kappa$, coefficients $a_j=\lambda(\tau)q_j$ and residual $(1-\lambda(\tau))\kappa$. It gives the asserted exact lower and upper volumes of $q_\tau$. The transverse recovery theorem applies, also confirming recovery directly.

For every $j$, the divisorial coefficient of $u_\tau$ along $H_j$ is positive. Projection preserves all ideals and hence these coefficients. Analytic singularities on a compact manifold have only finitely many divisorial components, since their defining coherent ideals have this property. Thus no analytic singularity type has the same all-weight ideals. The same conclusion holds upstairs along the distinct divisors $\pi^{-1}(H_j)$, by submersion pullback.
\end{proof}

\subsection{Exact pullback and smooth changes of background}

Let $\mu:Y\to X$ be a modification and let $\boldsymbol\Gamma$ satisfy \cref{thm:hopf-transverse-recovery}. By \cref{prop:i-envelope-modification}, its recovery identity pulls back to
\[
 \Gamma_\tau^\eta\circ\mu
       =P_{\mu^*(h+\eta)}[\Gamma_\tau\circ\mu]_{\mathcal I}.
\]
This concerns the exact pulled backgrounds. The form $\mu^*\eta$ may be degenerate, so these identities do not define a full system indexed by every Hermitian form on $Y$.

For a modification $\nu:\widetilde B\to B$, \cref{prop:singularity-volumes-modification} gives
\[
 \underline{\operatorname{vol}}_{\nu^*\kappa}
                  (\psi_a\circ\nu)
       =\underline{\operatorname{vol}}_\kappa(\psi_a)>0,
\]
so the transverse positivity test survives exact pullback. The same proposition also preserves ambient singularity volumes in an exact pulled background; it makes no comparison with an independent positive background upstairs.

Finally, the smooth transport $\mathsf G_f$ of \cref{prop:hermitian-system-gauge} commutes with both $\mathsf E$ and $\mathsf R$. Therefore recovery survives the change from $h$ to $h+\ddc f$ for every smooth real $f$. The transported scalar slices need not remain invariant under $E$, and no invariance of the raw masses of independently normalized nonclosed-background representatives is inferred. The general recovery problem without the scalar symmetry or transverse positive-volume condition remains open.
